\documentclass[a4paper,11pt]{article}

\usepackage[T2A]{fontenc}
\usepackage[utf8]{inputenc}
\usepackage[russian,english]{babel}
\usepackage[margin=2.3cm]{geometry}
\usepackage{parskip}
\usepackage{microtype}
\usepackage{xcolor}
\usepackage{listings}
\usepackage{hyperref}
\usepackage{booktabs}
\usepackage{longtable}
\usepackage{enumitem}
\usepackage{titlesec}
\usepackage{fancyhdr}
\usepackage{tikz}
\usetikzlibrary{arrows.meta,positioning,shapes.geometric,fit,backgrounds}

\hypersetup{colorlinks=true,linkcolor=blue!55!black,urlcolor=blue!55!black}

\lstdefinestyle{cpp}{
    language=C++,
    basicstyle=\ttfamily\footnotesize,
    keywordstyle=\color{blue!70!black}\bfseries,
    commentstyle=\color{green!45!black}\itshape,
    stringstyle=\color{red!60!black},
    numbers=left, numberstyle=\tiny\color{gray}, numbersep=6pt,
    frame=single, rulecolor=\color{gray!40},
    breaklines=true, showstringspaces=false, tabsize=4,
    captionpos=b,
}
\lstset{style=cpp}

\titleformat{\section}{\Large\bfseries}{\thesection}{0.8em}{}
\titleformat{\subsection}{\large\bfseries}{\thesubsection}{0.7em}{}
\titleformat{\subsubsection}{\normalsize\bfseries}{\thesubsubsection}{0.6em}{}

\pagestyle{fancy}
\fancyhf{}
\rhead{\small MirrorRoomSimulation}
\lhead{\small Документация разработчика}
\cfoot{\thepage}

\newcommand{\file}[1]{\texttt{#1}}
\newcommand{\code}[1]{\texttt{#1}}

\title{\textbf{MirrorRoomSimulation}\\[0.3em]
       \large Документация для программиста\\[0.2em]
       \normalsize (актуально для текущего состояния кода)}
\author{Репозиторий: \texttt{Githu228i/MirrorRoomSimulation}}
\date{\today}

\begin{document}
\maketitle

\begin{abstract}
\noindent
Этот документ — карта проекта для программиста, который открывает
\file{MirrorRoomSimulation} впервые. Здесь описано, что где лежит, как
устроена архитектура, где именно находится вычислительная логика
(спойлер: \emph{не} в \file{backendoperators.cpp} — там пусто),
как расширять код и какие в проекте есть подводные камни. Документ
написан по факту чтения исходников, а не по замыслу.
\end{abstract}

\tableofcontents
\newpage

% ============================================================
\section{О проекте}

\file{MirrorRoomSimulation} — учебное/демонстрационное приложение на C++
с GUI на Qt~6 (Widgets). Пользователь кликами мыши строит замкнутую
(или незамкнутую) цепочку зеркал в «комнате», редактирует точки и типы
зеркал через диалоги, а затем может трассировать луч, отражающийся от
зеркал.

\subsection{Стек}

\begin{longtable}{@{}p{4cm}p{10.5cm}@{}}
\toprule
\textbf{Компонент} & \textbf{Назначение} \\ \midrule
\endhead
C++17 & Язык. \\
Qt~6 (Widgets) & GUI, \code{QPainter}, диалоги, сигналы/слоты. \\
CMake $\geq$ 3.16 & Сборка. \\
Qt Creator & Рекомендуемая IDE. \\
\code{clangd} & Опционально. \\
\bottomrule
\end{longtable}

\subsection{Быстрый старт}

\begin{lstlisting}[language=bash,caption={Сборка}]
cmake -S . -B build/Desktop-Debug -DCMAKE_BUILD_TYPE=Debug
cmake --build build/Desktop-Debug -j
./build/Desktop-Debug/MirrorRoomSimulation
\end{lstlisting}

\begin{quote}
\textbf{Важно.} \file{build/}, \file{.qtcreator/}, \file{CMakeLists.txt.user},
\file{*.back*}, \file{.cache/} — не должны попадать в Git.
См.~\ref{sec:git}.
\end{quote}

% ============================================================
\section{Структура репозитория}

\begin{longtable}{@{}p{5.8cm}p{8.7cm}@{}}
\toprule
\textbf{Файл} & \textbf{Назначение} \\ \midrule
\endhead
\file{main.cpp} & \code{QApplication}, установка переводов, запуск \code{MainWindow}. \\
\file{mainwindow.h/.cpp} & Главное окно. Обработка мыши, \code{paintEvent}, вызов диалогов. \\
\file{mainwindow.ui} & Пустой дизайн-файл: только \code{centralwidget}, \code{menubar}, \code{statusbar}. Вся отрисовка — вручную в \code{paintEvent}. \\
\file{points.h/.cpp} & \code{Points} (базовая точка) и \code{Beam} (луч: точка + угол). \\
\file{mirror.h/.cpp} & Иерархия \code{Mirror} и её наследники. Пересечения зеркал. Рисование. \\
\file{polygon.h/.cpp} & \code{Polygon}: набор зеркал, замкнутость, валидация, трассировка луча. \\
\file{frontendoperators.h/.cpp} & Менеджер сцены: \code{CreatePoint}, \code{ConnectPoints}, \code{ChangeMirror}. \\
\file{backendoperators.h/.cpp} & \textbf{Пустая заглушка.} Только конструктор. Никакой логики нет. \\
\file{changeorconnectdialog.h/.cpp} & Диалог с двумя кнопками: «Изменить» / «Подключить». \\
\file{changepointdialog.h/.cpp} & Диалог редактирования точки. Определяет \code{ChangePointData}. \\
\file{MirrorRoomSimulation\_ru\_RU.ts} & Файл переводов (пустой). \\
\file{README.md} & Одна строка. \\
\file{CMakeLists.txt} & Сборка. \\
\bottomrule
\end{longtable}

\begin{quote}
\textbf{Ориентир.} Если вы ищете «где же логика backend» — идите в
\file{polygon.cpp} (трассировка луча) и \file{mirror.cpp} (пересечения).
\file{backendoperators.cpp} — это заготовка, которую, судя по всему,
планировали наполнить, но не наполнили.
\end{quote}

% ============================================================
\section{Архитектура}

\subsection{Слои}

\begin{enumerate}[label=\textbf{\arabic*.},leftmargin=*]
    \item \textbf{Геометрия-примитив:} \code{Points}, \code{Beam}.
    \item \textbf{Модель зеркал:} \code{Mirror} и наследники + геометрические
          алгоритмы пересечений (в анонимном namespace \file{mirror.cpp}).
    \item \textbf{Композиция сцены:} \code{Polygon} — держит вектор зеркал,
          трассирует луч, валидирует контур.
    \item \textbf{Менеджер сцены:} \code{FrontEndOperators} — обёртка над
          \code{Polygon} для UI-операций.
    \item \textbf{GUI:} \code{MainWindow} + диалоги.
\end{enumerate}

\begin{quote}
\textbf{Правило.} Модель (\code{Points}, \code{Mirror}, \code{Polygon}) не
должна знать о \code{QMainWindow}. Однако \code{Polygon} уже знает о
\code{QPainter} — наследники \code{Mirror} умеют себя рисовать. Это
осознанное упрощение (см.~\ref{sec:known-issues}).
\end{quote}

\subsection{Иерархия классов}

\begin{center}
\begin{tikzpicture}[
    every node/.style={draw, rounded corners, font=\footnotesize,
                       minimum height=0.7cm, align=center, inner sep=4pt},
    edge from parent/.style={draw, -Latex},
    level 1/.style={sibling distance=4cm, level distance=1.2cm},
    level 2/.style={sibling distance=3cm, level distance=1.2cm},
    level 3/.style={sibling distance=2.5cm, level distance=1.2cm},
    scale=0.95, transform shape
]
\node (Mirror) {Mirror\\\tiny (abstract)}
    child { node (Plane) {PlaneMirror} }
    child { node (NotPlane) {NotPlaneMirror}
        child { node (Conv) {ConvexMirror} }
        child { node (Conc) {ConcaveMirror} }
    };
\end{tikzpicture}
\end{center}

\begin{itemize}[leftmargin=*]
    \item \code{Mirror} — абстрактный базовый класс. Поля:
          \code{start\_}, \code{end\_}, \code{rad\_}.
          Виртуальный деструктор обязателен — хранится в
          \code{std::unique\_ptr<Mirror>}.
    \item \code{PlaneMirror} — плоское зеркало (rad = 0). Переопределяет
          \code{isFlat()} $\to$ \code{true}.
    \item \code{NotPlaneMirror} — промежуточный, защищённый конструктор.
    \item \code{ConvexMirror} — \code{rad > 0}, центр окружности «слева»
          от хорды.
    \item \code{ConcaveMirror} — \code{rad < 0}, центр «справа».
\end{itemize}

\subsection{Поток данных}

\begin{enumerate}[leftmargin=*]
    \item Пользователь кликает в \code{MainWindow::mousePressEvent}.
    \item \code{MainWindow} решает: попадание в точку, близко, новая точка.
    \item Вызывается один из методов \code{FrontEndOperators}:
          \code{CreatePoint}, \code{ConnectPoints} или \code{ChangeMirror}.
    \item Тот манипулирует \code{Polygon::sides\_}
          (\code{std::vector<std::unique\_ptr<Mirror>>}).
    \item \code{MainWindow} вызывает \code{update()} — перерисовка.
    \item \code{MainWindow::paintEvent} рисует \code{first} и все зеркала
          через их виртуальный \code{Draw()}.
\end{enumerate}

% ============================================================
\section{Ключевые сущности}

\subsection{Points и Beam (\file{points.h})}

\begin{lstlisting}[caption={Points / Beam}]
class Points {
public:
    Points(double x = 0.0, double y = 0.0);
    virtual ~Points() = default;
    double getX() const;
    double getY() const;
    double Dist(const Points& other) const;
    virtual std::string Cords() const;
protected:
    double x_, y_;
};

class Beam : public Points {
public:
    Beam(double x = 0.0, double y = 0.0, double angle = 0.0);
    double getA() const;
    std::string Cords() const override;
private:
    double angle_;
};
\end{lstlisting}

\begin{quote}
\code{Cords()} возвращает строку формата \code{"x y"} (или
\code{"x y angle"} для луча) с точностью 3 знака. Используется только
для отладочной записи в файл.
\end{quote}

\subsection{Иерархия Mirror (\file{mirror.h})}

Ключевые методы \code{Mirror}:

\begin{longtable}{@{}p{6.3cm}p{8.2cm}@{}}
\toprule
\textbf{Метод} & \textbf{Назначение} \\ \midrule
\endhead
\code{Mirror(st, end, rad)} & Бросает \code{std::invalid\_argument} при нулевой хорде или $|rad| < d/2$. \\
\code{isFlat()} & \code{true}, если $|rad| < 10^{-9}$ или $> 10^{7}$. \\
\code{CircleCenter()} & Центр окружности; сторона задаётся знаком \code{rad}. \\
\code{CircleCenter(bool leftSide)} & Явное задание стороны. \\
\code{PointOnArc(p, c)} & Лежит ли точка на дуге/отрезке. \\
\code{Intersections(other)} & Вектор точек пересечения. \\
\code{Touches(other)} & \code{!Intersections().empty()}. \\
\code{Crossing(other)} & Пересекаются не только в общих концах. \\
\code{Draw(painter, insideOnLeft)} & Виртуальный. Чисто виртуальный в базе. \\
\bottomrule
\end{longtable}

Знак \code{rad\_}:
\begin{itemize}[leftmargin=*]
    \item $rad > 0$ — центр \emph{слева} от направления \code{start → end}.
          Соответствует \code{ConvexMirror}.
    \item $rad < 0$ — центр \emph{справа}. Соответствует \code{ConcaveMirror}.
    \item $rad = 0$ — отрезок (\code{PlaneMirror}).
\end{itemize}

\subsection{Polygon (\file{polygon.h})}

Хранилище сцены.

\begin{lstlisting}[caption={Поля Polygon}]
class Polygon {
    // ...
private:
    std::vector<std::unique_ptr<Mirror>> sides_;
    Points first;             // первая точка незамкнутого контура
    bool isExist_ = false;    // создана ли первая точка
    bool isCycled_ = false;   // замкнут ли контур
    Beam beam_;               // луч для трассировки
};
\end{lstlisting}

Ключевые методы:

\begin{longtable}{@{}p{5.5cm}p{9cm}@{}}
\toprule
\textbf{Метод} & \textbf{Что делает} \\ \midrule
\endhead
\code{Push(const Mirror\&)} & Копирует зеркало через \code{dynamic\_cast} и добавляет в \code{sides\_}. Вызывает \code{Write()}. \\
\code{Del(int idx)} & Удаляет зеркало по индексу. Вызывает \code{Write()}. \\
\code{Check()} & Валидация: контур замкнут, стороны пересекаются только в общих вершинах соседей. \\
\code{TraceBeam(maxBounces)} & Трассировка луча. Пишет \code{beam.txt}. \\
\code{Write()} & Дамп вершин и луча в \code{points.txt}. \\
\bottomrule
\end{longtable}

\begin{quote}
\textbf{Побочный эффект.} \code{Push}, \code{Del}, \code{Write},
\code{TraceBeam} открывают файлы. При интенсивном редактировании это
может тормозить и засорять рабочую директорию.
\end{quote}

\subsection{FrontEndOperators (\file{frontendoperators.h})}

\begin{lstlisting}[caption={FrontEndOperators}]
class FrontEndOperators {
public:
    FrontEndOperators();
    void CreatePoint(double x, double y);
    void ConnectPoints(int num);
    void ChangeMirror(int num, const ChangePointData &data);
    friend class MainWindow;
private:
    Polygon prototype;
};
\end{lstlisting}

\begin{quote}
\textbf{Имя вводит в заблуждение.} Это \emph{менеджер сцены}, а не слой
отрисовки. Отрисовка — в \code{MainWindow::paintEvent} и
\code{Mirror::Draw}.
\end{quote}

Логика методов:
\begin{itemize}[leftmargin=*]
    \item \code{CreatePoint} — добавляет новую точку.
          Если \code{isExist\_} ложно — сохраняет \code{first} и ставит
          флаг. Иначе строит \code{PlaneMirror} от конца последнего
          зеркала к новой точке. Лимит: 8 точек
          (\code{if (prototype.sides\_.size() < 8)}).
    \item \code{ConnectPoints(num)} — замыкает контур: удаляет
          зеркала до \code{num}, добавляет последнее зеркало от
          \code{end} к \code{start}, ставит \code{isCycled\_ = true}.
    \item \code{ChangeMirror(num, data)} — самая сложная функция.
          Разбирает случаи: последняя точка незамкнутого контура,
          первая точка замкнутого, обычная точка. Везде через
          \code{dynamic\_cast} определяет тип зеркала и пересоздаёт
          соседние.
\end{itemize}

\subsection{MainWindow (\file{mainwindow.h})}

Поля всего два:
\begin{lstlisting}
Ui::MainWindow *ui;
FrontEndOperators polygonCreator;
\end{lstlisting}

Логика:
\begin{itemize}[leftmargin=*]
    \item \code{mousePressEvent} — вся обработка кликов. Хардкод:
          \code{y > 1000} отсекается, радиус попадания в точку
          \code{25}, минимальное расстояние между точками \code{100}.
    \item \code{paintEvent} — рисует \code{first}, линию от \code{first}
          к \code{temp[0]->getStart()} и все зеркала.
\end{itemize}

\subsection{Диалоги}

\code{ChangeOrConnectDialog} — простое окно с двумя кнопками
(\code{accept} / \code{reject}). Текст-заглушка:
\code{"ТУТ БУДЕТ ТВОЙ ТЕКСТ"}.

\code{ChangePointDialog} — редактирование одной точки. Возвращает
структуру:

\begin{lstlisting}
struct ChangePointData {
    double x;
    double y;
    int mirrorType;   // 0 = plane, 1 = concave, 2 = convex
    double radius;
};
\end{lstlisting}

Если \code{lastPoint == true} — комбобокс типа зеркала и поле радиуса
отключены (у последней точки зеркала ещё нет).

\subsection{BackEndOperators — пустая заглушка}

\begin{lstlisting}[caption={Весь файл backendoperators.cpp}]
#include "backendoperators.h"
BackEndOperators::BackEndOperators() {}
\end{lstlisting}

\begin{quote}
\textbf{Важно.} Не ищите здесь логику. Её нет. Вся вычислительная
часть живёт в \file{mirror.cpp} и \file{polygon.cpp} в безымянных
\code{namespace \{ \}}. Если вы хотите что-то добавить «в backend» —
логичнее либо дописать \code{BackEndOperators}, либо создать новый
модуль, а не размазывать по \file{polygon.cpp}.
\end{quote}

% ============================================================
\section{Взаимодействие пользователя}

\subsection{Создание точки}

Клик в пустое место:
\begin{enumerate}[leftmargin=*]
    \item Если попадание в существующую точку ($\le 25$) — редактирование.
    \item Если близко к существующей ($\le 100$) — сообщение «Точки слишком близко!».
    \item Иначе — \code{CreatePoint(x, y)}.
\end{enumerate}

Лимит: \textbf{8 сторон}.

\subsection{Редактирование точки}

Клик по существующей точке:
\begin{itemize}[leftmargin=*]
    \item \textbf{Контур замкнут} — сразу открывается
          \code{ChangePointDialog}.
    \item \textbf{Контур не замкнут} — сначала
          \code{ChangeOrConnectDialog} («Изменить» или «Подключить»).
\end{itemize}

\subsection{Соединение точек}

Выбор «Подключить» в незамкнутом контуре:
\begin{itemize}[leftmargin=*]
    \item Последнее зеркало подключить нельзя —
          \code{QMessageBox::critical}.
    \item Иначе \code{ConnectPoints(i)}: удаляет все зеркала до
          \code{i}, добавляет одно замыкающее.
\end{itemize}

% ============================================================
\section{Геометрические алгоритмы}

\subsection{Пересечения (\file{mirror.cpp}, анонимный namespace)}

\begin{itemize}[leftmargin=*]
    \item \code{IntersectSegments} — отрезок/отрезок,
          включая коллинеарный случай с вычислением перекрытия.
    \item \code{IntersectSegmentArc} — отрезок/дуга, квадратное
          уравнение.
    \item \code{IntersectArcs} — дуга/дуга, стандартная формула
          пересечения окружностей с проверкой \code{PointOnArc}.
\end{itemize}

Константы:
\code{EPS = 1e-9}, \code{POINT\_EPS = 1e-6}, \code{PAR\_EPS = 1e-10},
\code{FLAT\_R = 1e7}.

\subsection{Трассировка луча (\file{polygon.cpp})}

\code{TraceBeam}:
\begin{enumerate}[leftmargin=*]
    \item Позиция и направление из \code{beam\_}.
    \item Для каждого зеркала — пересечение луча.
    \item Выбор ближайшего с $t > 10^{-7}$.
    \item Нормаль в точке удара (\code{GetNormal}).
    \item Отражение \code{d' = d - 2(d·n)n}
          (\code{ReflectDirection}), с перенормировкой.
    \item Запись в \code{beam.txt}: позиции через пробел.
\end{enumerate}

% ============================================================
\section{Файловый ввод/вывод}

\begin{longtable}{@{}p{4cm}p{4.5cm}p{6cm}@{}}
\toprule
\textbf{Файл} & \textbf{Кто пишет} & \textbf{Формат} \\ \midrule
\endhead
\file{points.txt} & \code{Polygon::Write()} & Заголовок, координаты вершин, луч. \\
\file{beam.txt} & \code{Polygon::TraceBeam()} & x y для каждой точки удара. \\
\bottomrule
\end{longtable}

\begin{quote}
\textbf{Внимание.} Оба файла пишутся в текущую рабочую директорию
(не в \code{\$HOME} и не в директорию проекта). При запуске из Qt Creator
это \code{build/Desktop-Debug/}. Легко закоммитить по ошибке — держите
их в \code{.gitignore}.
\end{quote}

% ============================================================
\section{Сборка и CMake}

\begin{lstlisting}[caption={CMakeLists.txt, фрагмент}]
set(PROJECT_SOURCES
    main.cpp
    mainwindow.cpp mainwindow.h mainwindow.ui
    mirror.h mirror.cpp
    points.h points.cpp
    frontendoperators.h frontendoperators.cpp
    backendoperators.h backendoperators.cpp
    polygon.h polygon.cpp
)
qt_add_executable(MirrorRoomSimulation
    ${PROJECT_SOURCES}
    changeorconnectdialog.h changeorconnectdialog.cpp
    changepointdialog.h changepointdialog.cpp
)
\end{lstlisting}

\begin{quote}
\textbf{Странность.} Диалоги добавлены отдельным блоком в
\code{qt\_add\_executable}, а не в \code{PROJECT\_SOURCES}. Если
добавляете новый файл — выбирайте одно место и следуйте стилю.
\end{quote}

\subsection{Переводы}

\begin{lstlisting}[language=bash]
lupdate . -ts MirrorRoomSimulation_ru_RU.ts
lrelease MirrorRoomSimulation_ru_RU.ts
\end{lstlisting}

\begin{quote}
\textbf{Баг.} В \file{main.cpp} имя перевода —
\code{"MirrorRoomSimulationProject\_" + locale}, а CMake генерирует
\code{MirrorRoomSimulation\_}. Префикс \code{Project} лишний.
Либо приведите \file{main.cpp} к \code{"MirrorRoomSimulation\_"},
либо переименуйте \code{TS\_FILES}. Плюс сам \file{.ts} пустой —
переводить пока нечего.
\end{quote}

% ============================================================
\section{Точки расширения}
\label{sec:extending}

\begin{longtable}{@{}p{6cm}p{9cm}@{}}
\toprule
\textbf{Задача} & \textbf{Куда идти} \\ \midrule
\endhead
Добавить новый тип зеркала &
Наследник \code{Mirror} + \code{Draw}. Плюс: конструктор в
\code{Polygon::Push}, \code{ChangeMirror}, \code{dynamic\_cast}-ветки. \\
\addlinespace
Изменить физику отражения &
\code{ReflectDirection} и \code{GetNormal} в \file{polygon.cpp}
(анонимный namespace). \\
\addlinespace
Добавить тип геометрического объекта &
\code{Points} или отдельный класс. \\
\addlinespace
Изменить отрисовку &
\code{Mirror::Draw} в наследниках, \code{drawCircularArc} в
\file{mirror.cpp}, \code{MainWindow::paintEvent}. \\
\addlinespace
Добавить пункт меню / горячую клавишу &
\code{mainwindow.ui} (пустой) или \code{MainWindow} в коде. \\
\addlinespace
Снять лимит в 8 точек &
\code{FrontEndOperators::CreatePoint}, условие
\code{sides\_.size() < 8}. \\
\addlinespace
Настроить трассировку (число отражений) &
Параметр \code{maxBounces} у \code{TraceBeam}. Сейчас не вызывается
из UI. \\
\bottomrule
\end{longtable}

\subsection{Чек-лист при добавлении файла}

\begin{enumerate}[leftmargin=*]
    \item Создать \file{.h} и \file{.cpp}.
    \item Добавить в \code{CMakeLists.txt} (в \code{PROJECT\_SOURCES}
          или рядом с диалогами — но не в оба места сразу).
    \item Если класс — \code{QObject}, добавить \code{Q\_OBJECT}.
    \item Если есть переводимые строки — \code{tr(...)}.
    \item Собрать, убедиться, что нет предупреждений moc/uic.
    \item Закоммитить только исходники и CMake, без \file{build/}.
\end{enumerate}

% ============================================================
\section{Известные проблемы и странности}
\label{sec:known-issues}

\begin{longtable}{@{}p{6.5cm}p{8cm}@{}}
\toprule
\textbf{Что} & \textbf{Где / почему} \\ \midrule
\endhead
\code{BackEndOperators} пуст &
\file{backendoperators.cpp}. Логика размазана по \file{polygon.cpp}
и \file{mirror.cpp}. \\
\addlinespace
Несоответствие имён переводов &
\file{main.cpp} ищет \code{MirrorRoomSimulationProject\_*},
CMake генерирует \code{MirrorRoomSimulation\_*}. \\
\addlinespace
Магические числа &
\code{1600}, \code{1200}, \code{1000}, \code{25}, \code{100},
\code{8}, \code{4}. Все в \file{mainwindow.cpp} и
\file{frontendoperators.cpp}. \\
\addlinespace
\code{Polygon::Write()} на каждый \code{Push}/\code{Del} &
Открывает \file{points.txt} при каждом изменении. \\
\addlinespace
\code{dynamic\_cast}-цепочки &
\code{Polygon::Push}, \code{FrontEndOperators::ChangeMirror}. При
добавлении нового типа зеркала придётся дописывать везде. \\
\addlinespace
\code{MainWindow} — friend & У \code{Polygon} и
\code{FrontEndOperators}. Ломает инкапсуляцию. \\
\addlinespace
\code{mirror.cpp} включает \code{QPainter} &
Модель знает о GUI. Смешение слоёв. \\
\addlinespace
Текст-заглушка в диалоге &
\code{"ТУТ БУДЕТ ТВОЙ ТЕКСТ"} в
\file{changeorconnectdialog.cpp}. \\
\addlinespace
\file{README.md} пуст & Одна строка. \\
\addlinespace
\code{ChangePointData} в \file{changepointdialog.h} &
\code{FrontEndOperators} тянет диалоговый заголовок только ради
структуры. Логичнее вынести в отдельный \file{changepointdata.h}. \\
\bottomrule
\end{longtable}

% ============================================================
\section{Правила работы с Git}
\label{sec:git}

\subsection{Что не коммитить}

\begin{itemize}[leftmargin=*]
    \item \file{build/}, \file{build-*/}, \file{cmake-build-*/};
    \item \file{CMakeLists.txt.user}, \file{*.user}, \file{.qtcreator/};
    \item \file{.cache/}, \file{compile\_commands.json};
    \item \file{*.back}, \file{*.back1}, \file{*.back2}, \file{*~},
          \file{*.orig}, \file{*.rej};
    \item \file{points.txt}, \file{beam.txt} (артефакты работы).
\end{itemize}

\begin{lstlisting}[language=bash,caption={.gitignore}]
build/
build-*/
cmake-build-*/
CMakeLists.txt.user
*.user
*.user.*
.qtcreator/
.cache/
compile_commands.json
*.back
*.back1
*.back2
*.back3
*~
*.orig
*.rej
points.txt
beam.txt
\end{lstlisting}

\subsection{Если мусор уже в индексе}

\begin{lstlisting}[language=bash]
git rm -r --cached build/ .qtcreator/ CMakeLists.txt.user
git rm --cached points.txt beam.txt
git commit -m "Stop tracking build artifacts and IDE files"
\end{lstlisting}

\subsection{Стратегия веток}

Ветки \code{main} и \code{backend} имеют \emph{разную структуру файлов}.
\code{git merge} между ними даёт десятки конфликтов типа
«modify/delete» и не имеет смысла.

\begin{itemize}[leftmargin=*]
    \item Если \code{backend} — актуальная версия:
\begin{lstlisting}[language=bash]
git branch main-backup main     # резервная копия
git checkout backend
git branch -f main backend
git checkout main
git push --force-with-lease origin main
\end{lstlisting}
    \item Если нужно перенести отдельные коммиты —
          \code{git cherry-pick <hash>}.
    \item Избегайте \code{git merge} между ветками с разной структурой
          файлов: Git не сопоставляет \file{mirror.cpp} и
          \file{backendoperators.cpp} как один модуль.
\end{itemize}

% ============================================================
\section{Соглашения по коду}

\begin{itemize}[leftmargin=*]
    \item Классы — \code{CamelCase}; методы — \code{camelCase};
          приватные поля — \code{camelCase\_} (например, \code{sides\_},
          \code{rad\_}).
    \item Заголовки — include guards, не \code{\#pragma once}
          (в проекте везде guards).
    \item Не использовать \code{using namespace std;}.
    \item Forward-declaration в \file{.h} (\code{class QPainter;})
          вместо \code{\#include <QPainter>} — как в \file{mirror.h}.
    \item Публичные методы — с кратким комментарием о
          предусловиях (как у \code{Mirror}).
    \item Новые типы зеркал — наследовать от \code{Mirror}, реализовать
          чисто виртуальный \code{Draw}.
\end{itemize}

% ============================================================
\section{Отладка и типовые проблемы}

\begin{longtable}{@{}p{6.3cm}p{8.2cm}@{}}
\toprule
\textbf{Симптом} & \textbf{Что делать} \\ \midrule
\endhead
\code{undefined reference to vtable} &
Забыт \code{Q\_OBJECT} (для QObject-классов) или файл не добавлен
в CMake. \\
\addlinespace
\code{moc\_*.cpp} не найден &
Проверить \code{AUTOMOC ON}. \\
\addlinespace
Не рисуется новая точка &
Не вызван \code{update()} после изменения модели. \\
\addlinespace
Сцена не перерисовывается после \code{ChangeMirror} &
Проверить возврат из \code{ChangeMirror} — там \code{return} в нескольких
ветках. \\
\addlinespace
\code{std::invalid\_argument} из \code{Mirror} &
Радиус меньше половины хорды или нулевая хорда. \\
\addlinespace
Переводы не подхватываются &
Префикс \code{MirrorRoomSimulationProject\_} в \file{main.cpp}
не совпадает с CMake. См.~раздел «Переводы». \\
\addlinespace
В рабочей папке появляется \file{points.txt} &
Это \code{Polygon::Write()}. Добавить в \code{.gitignore}. \\
\addlinespace
Qt Creator создаёт \file{*.back1/.back2} &
Добавить в \code{.gitignore}, удалить \code{rm -f *.back*}. \\
\bottomrule
\end{longtable}

% ============================================================
\section{Шпаргалка}

\begin{lstlisting}[language=bash,caption={Часто используемое}]
# Сборка
cmake --build build/Desktop-Debug -j

# Пересборка с нуля
rm -rf build/Desktop-Debug && \
  cmake -S . -B build/Desktop-Debug && \
  cmake --build build/Desktop-Debug -j

# Очистка от мусора Qt Creator и рабочих файлов
rm -f *.back *.back1 *.back2 *.back3 *~ points.txt beam.txt

# Проверка, что мусор не в индексе
git ls-files | grep -E '^(build/|\.qtcreator/)' && echo "FOUND" || echo "clean"

# Обновление переводов
lupdate . -ts MirrorRoomSimulation_ru_RU.ts
lrelease MirrorRoomSimulation_ru_RU.ts
\end{lstlisting}

% ============================================================
\section{Куда смотреть в первую очередь}

\begin{enumerate}[leftmargin=*]
    \item \file{main.cpp} — запуск.
    \item \file{mainwindow.h} — состояние окна.
    \item \file{mainwindow.cpp} — \code{mousePressEvent} (главный вход
          пользователя) и \code{paintEvent} (отрисовка).
    \item \file{frontendoperators.h} — операции над сценой.
    \item \file{frontendoperators.cpp} — \code{CreatePoint},
          \code{ConnectPoints}, \code{ChangeMirror}.
    \item \file{polygon.h} — модель сцены.
    \item \file{polygon.cpp} — \code{TraceBeam}, отражение.
    \item \file{mirror.h} — иерархия зеркал.
    \item \file{mirror.cpp} — пересечения и \code{Draw}.
    \item Диалоги — по мере необходимости.
\end{enumerate}

\vspace{1em}
\hrule
\vspace{0.4em}
\noindent\small\textit{Документ написан по факту чтения исходников.
При изменении структуры обновляйте разделы 2, 3, 9, 10.}

\end{document}
